\documentclass[12pt,a4paper]{article}
\usepackage[utf8]{inputenc}
\usepackage[T2A]{fontenc}
\usepackage[russian]{babel}
\usepackage{amsmath,amssymb}
\usepackage{geometry}
\geometry{margin=2.3cm}
\DeclareMathOperator{\Tr}{Tr}
\title{Граничное роторное действие\\для вильсоновского функционала}
\author{}
\date{6 августа 2026 г.}
\begin{document}
\maketitle

\section{Целевой функционал}

Требуется получить из устойчивого локального механизма
\begin{equation}
 V_{\mathrm{gap}}(c)
 =\frac8{1-c}+\frac83\log(1-c)-\frac{44}{45}c.
\end{equation}
Обычное исключение поля с фиксированным источником даёт неверный
отрицательный знак. Положительный обратный член возникает у циклической
степени свободы в ансамбле фиксированного заряда:
\begin{equation}
 H_{\mathrm{rot}}=\frac{q^2}{2m^2},
 \qquad m^2(c)=1-c.
\end{equation}

\section{Минимальные кратности}

Для \(N\) вещественных роторов, нормированного следа ранга \(d\) и средней
нормы заряда \(q^2\)
\begin{equation}
 \Gamma_{\det}=\frac{N}{2d}\log(1-c),
 \qquad
 \Gamma_{\mathrm{rot}}=\frac{Nq^2}{2d(1-c)}.
\end{equation}
Точные коэффициенты требуют
\begin{equation}
 \frac{N}{2d}=\frac83,
 \qquad
 \frac{Nq^2}{2d}=8.
\end{equation}
Минимальное целочисленное решение:
\begin{equation}
 N=16,\qquad d=3,\qquad q^2=3.
\end{equation}

\section{Каноническое представление}

Единственное приводимое представление размерности \(16\), не содержащее
синглетов и состоящее только из полуцелых спинов, со средним казимиром \(3\):
\begin{equation}
 \boxed{\mathcal R_{\partial}=2V_{1/2}\oplus3V_{3/2}.}
\end{equation}
Для него
\begin{equation}
 \dim\mathcal R_{\partial}=16,
 \qquad
 \Tr_{\mathcal R_{\partial}}C_2=48.
\end{equation}
При \(d=3\)
\begin{equation}
 \frac{\dim\mathcal R_{\partial}}{2d}=\frac83,
 \qquad
 \frac{\Tr_{\mathcal R_{\partial}}C_2}{2d}=8.
\end{equation}
Все блоки имеют полуцелый спин, поэтому центр \(SU(2)\) действует знаком
минус и одновременно задаёт антипериодическую ветвь.

\section{Открытые условия}

Необходимо вывести кратности \(2\) и \(3\), явный граничный оператор,
линейный член \(-44c/45\), полный комплекс БВ/БРСТ и связь с осью поколений.
До этого конструкция является конкретным локальным кандидатом, но не
завершённым действием.

\end{document}